% Drawing 156-6/04-1 -- folder 156-6, batch 04, archive page 65 (PDF page 66).
% Transcription: schéma : $F' \ll F$ sur la première ligne ; sous $F'$, $X' \leq F'$ ; sous $F$, $G \leq F$ ; sous $X'$, $Z \overset{\circ}{\ll} X'$ (signe pointé d'un petit cercle), et $Z \ll G$ sur la dernière ligne
% Kind: order diagram F' << F, X' <= F', G <= F, Z <<o X', Z << G. Method 2 (schematic redraw from the facsimile). Status: ok.
% Notes: Relation diagram written with symbols: F' << F on the first line; under F', X' <= F' (the <= turned upright); under F, G <= F; under X', Z <<o X' (the << turned, marked with a small circle); Z << G on the last line (the << slightly slanted).
% Coordinates: 1 unit = 1 mm = 1/100 of the crop width (bbox in page fractions: [0.48, 0.42, 0.75, 0.54]).
% Redrawn by Claude (claude-opus-5-5), issue #73, 2026-11-07; not checked by a human; not compiled (no TeX in the session).
% Rights: reproduces part of the work -- behind RENDER_TRANSCRIPTS, never published by itself.
\documentclass[tikz,border=3pt]{standalone}
\usepackage{amsmath,amssymb}
\usetikzlibrary{patterns,arrows.meta}
\tikzset{pencil/.style={black!45}, struck/.style={black!35}, mark/.style={circle,draw,inner sep=1.3pt}, pt/.style={circle,fill,inner sep=0.9pt}, opt/.style={circle,draw,fill=white,inner sep=1.1pt}}
\begin{document}
\begin{tikzpicture}[x=1mm,y=1mm,line cap=round,line join=round,font=\scriptsize]
\node[inner sep=1pt] at (31,-13.6) {\scriptsize $F'$};
\node[inner sep=1pt] at (53,-12.3) {\scriptsize $\ll$};
\node[inner sep=1pt] at (68,-10.4) {\scriptsize $F$};
\node[inner sep=1pt,rotate=90] at (31,-21.4) {\scriptsize $\leq$};
\node[inner sep=1pt,rotate=90] at (70,-24.7) {\scriptsize $\leq$};
\node[inner sep=1pt] at (31,-29.2) {\scriptsize $X'$};
\node[inner sep=1pt] at (72,-39) {\scriptsize $G$};
\node[inner sep=1pt,rotate=90] at (28,-40.3) {\scriptsize $\ll$};
\node[inner sep=1pt] at (21,-39.6) {\tiny $\circ$};
\node[inner sep=1pt] at (25.3,-50) {\scriptsize $Z$};
\node[inner sep=1pt,rotate=-15] at (50,-41.6) {\scriptsize $\ll$};
\end{tikzpicture}
\end{document}
